\chapter{Enveloping bars and planar path sequences}
\label{ch:enveloping-audit}

The dimension of a bar chain space and the dimension of its homology are different invariants. For an enveloping algebra, the Chevalley comparison gives a Lie-theoretic model for the homology. This makes it possible to compare Motzkin and Riordan path counts with the ordered bar on a specified carrier.

Bar length and Chevalley degree are homological degrees: $H_n$ denotes degree-$n$ homology. The cohomology of the dual Chevalley complex is written $H^n$. For finite-dimensional weight pieces their dimensions agree.

\section{Chevalley comparison}

\begin{theorem}[Bar--Chevalley comparison]
\label{thm:bar-chevalley-spine}
Let $k$ be a field of characteristic zero and let $\g$ be an ordinary Lie algebra over $k$.
Give $U(\g)$ its augmentation killing $\g$.
There is a natural quasi-isomorphism of ordinary conilpotent differential graded coalgebras
\[
 \kappa:C_*(\g)\longrightarrow\BarPerp\bigl(U(\g)\bigr),
\]
where both constructions use direct sums and ordinary tensors over $k$.
It preserves the PBW filtration.
If $\g=\bigoplus_{w\ge1}\g_w$ and $[\g_u,\g_v]\subseteq\g_{u+v}$,
the map is a quasi-isomorphism in each total weight.
It extends to the degreewise product over weights, using completed tensor products by total weight.
The completed coalgebras are conilpotent in each weight; no global conilpotence is asserted for infinite sums.
\end{theorem}
\begin{proof}
Put $C_n(\g)=\Lambda^n\g$ in homological degree $n$, with differential
\[
 d(x_1\wedge\cdots\wedge x_n)
 =\sum_{i<j}(-1)^{i+j+1}[x_i,x_j]\wedge
 x_1\wedge\cdots\widehat{x_i}\cdots\widehat{x_j}\cdots\wedge x_n.
\]
Its coproduct splits the exterior factors between two ordered sublists,
with the sign of the corresponding permutation.
The ordinary normalized bar differential joins positions $i,i+1$ with sign $(-1)^{i-1}$.
Define
\[
 \kappa(x_1\wedge\cdots\wedge x_n)
 =\sum_{\sigma\in S_n}\operatorname{sgn}(\sigma)
 [x_{\sigma(1)}|\cdots|x_{\sigma(n)}].
\]
At degree zero this map preserves the unit.
At degree two its differential is $[x_1x_2-x_2x_1]=[[x_1,x_2]]$.
In every degree, pair permutations whose merged adjacent entries occur in opposite orders.
Their difference is the bracket of those entries, with the displayed exterior sign.
This proves the chain-map equation.
Cutting an antisymmetrized word amounts to choosing an ordered subset and
antisymmetrizing its two subwords.
The two procedures give the same permutations and signs, proving the coalgebra identity.

Filter $U(\g)$ by PBW degree, the bar by total PBW degree, and $\Lambda^n\g$ by degree $n$.
PBW identifies the associated graded algebra with $S(\g)$.
The associated graded comparison is the antisymmetrization
$\Lambda(\g)\to B(S(\g))$, with zero differential on its source.
To verify that it is a quasi-isomorphism, consider the free Koszul resolution
$S(V)\otimes\Lambda(V)$ for $V=\g$.
For a basis $(x_i)$ of $V$, its differential is
$\delta=\sum_i x_i\iota_i$, where $\iota_i$ contracts the exterior factor $x_i$.
On polynomial degree $p$ and exterior degree $q$, with $p+q>0$, the map
\[
 h=\frac{1}{p+q}\sum_i (x_i\wedge -)\,\frac{\partial}{\partial x_i}
\]
satisfies $\delta h+h\delta=1$.
Only finitely many summands act on each element, even when $V$ is infinite-dimensional.
The identity follows by cancellation of cross terms and the Euler identity in the two degrees.
Thus this free resolution has augmentation homology $k$ only.
The free normalized bar resolution has the same homology by its extra-degeneracy contraction.
Antisymmetrization defines an augmentation-preserving comparison between these free resolutions.
The simplicial bar sign convention is related to the displayed bar convention by multiplication by $(-1)^n$ in degree $n$.
Free bases permit a reverse comparison and homotopies to be lifted inductively,
since every positive-degree cycle in a resolution is a boundary.
Tensoring these homotopies with $k$ proves the asserted quasi-isomorphism on associated graded complexes.

Every cycle in the cone of $\kappa$ has finite PBW filtration degree.
A primitive for its highest associated graded part lowers that degree;
iteration terminates because the filtration is bounded below.
The cone is acyclic, proving the ordinary assertion.
For a positive weight grading, PBW degree is at most total weight.
The same argument is therefore finite within each weight.
Finally, kernels and images of a degreewise product differential are products of kernels and images,
so the componentwise quasi-isomorphisms extend to the stated completion.
\end{proof}

The ordinary comparison does not itself construct a morphism in $\FactD(X)$.
For geometric coefficient systems, one must give compatible enveloping and Chevalley constructions on every chart
and coherent factorization equivalences for the whole Chevalley complex.
On each disjointness locus, with restriction denoted by $j^!$, the required equation is
\[
 f_{B,I}\,j^!\kappa_I
 =(\kappa_{I_1}\boxtimes\kappa_{I_2})\,f_{C,I},
 \qquad I=I_1\sqcup I_2.
\]
Here $f_{C,I}$ and $f_{B,I}$ are the specified Chevalley and bar factorization maps.
The bar map may use the signed shuffle comparison under its geometric functor hypotheses.
The Chevalley factorization maps, collision compatibility, refinement coherence,
and continuity of the completed maps remain additional geometric data to construct or verify.
Internal tensor formulas and PBW alone do not provide them.


\begin{calculation}[Finite-dimensional simple case]
\label{calc:simple-spine}
For a finite-dimensional complex simple Lie algebra of rank $r$ with
exponents $m_1,\dots,m_r$, the
Chevalley cohomology is the exterior algebra on primitive classes of degrees
$2m_i+1$. The corresponding homology Poincar\'e polynomials are
\[
 \BarPerp U(\mathfrak{sl}_2):\ 1+t^3,
 \qquad
 \BarPerp U(\mathfrak{sl}_3):\ 1+t^3+t^5+t^8,
 \qquad
 \BarPerp U(\mathfrak g_2):\ 1+t^3+t^{11}+t^{14}.
\]
The corresponding Betti sequences for $\mathfrak{sl}_2$ and $\mathfrak{sl}_3$ are
$1,0,0,1$ and $1,0,0,1,0,1,0,0,1$.
For a finite-dimensional Lie algebra, the Chevalley chain dimensions are
$\binom{\dim\g}{n}$.

For an ordinary augmented algebra $A$ over $\mathbb C$, let
$I=\ker(A\to\mathbb C)$.  Its normalized two-sided bar with trivial
coefficients has homological chain spaces
\[
 B_0(A)=\mathbb C,\qquad B_n(A)=I^{\otimes n}\quad(n\ge1).
\]
Here $n$ is bar length, independently of any internal weight or PBW degree.
If $I$ has finite dimension $d$, then $\dim B_n(A)=d^n$; in particular,
an augmentation ideal of dimension three gives the chain counts $3^n$.
In a cohomological convention that places $B_n$ in degree $-n$, its homology
is denoted $H^{-n}$ instead of $H_n$.

For $A=U(\mathfrak{sl}_2)$, PBW gives the basis
$e^a h^b f^c$, with $a,b,c\ge0$, where
$[h,e]=2e$, $[h,f]=-2f$, and $[e,f]=h$.
The augmentation ideal has all these basis elements except $1$.
Consequently $B_n(A)$ is countably infinite-dimensional for every $n\ge1$.
Its total PBW filtration has finite-dimensional associated graded pieces,
with generating function
\[
 \sum_{d\ge0}\dim\operatorname{gr}^{\mathrm{PBW}}_d B_n(A)\,z^d
 =\bigl((1-z)^{-3}-1\bigr)^n.
\]
In particular, $3^n$ is the dimension in the least possible total PBW degree
$d=n$. These least-degree dimensions are $3,9,27,81,243$ for
$n=1,2,3,4,5$; the full bar chain spaces are infinite-dimensional. The quotient of
indecomposables is zero: the perfectness of $\mathfrak{sl}_2$ implies $I/I^2=0$.

The Chevalley chain model is different: its degree-$n$ space is
$\Lambda^n\mathfrak{sl}_2$, of dimensions $1,3,3,1$ in degrees $0,1,2,3$.
With the convention $d(x\wedge y)=[x,y]$, the map
$d:\Lambda^2\mathfrak{sl}_2\to\mathfrak{sl}_2$ sends
\[
 e\wedge h\longmapsto-2e,\qquad
 e\wedge f\longmapsto h,\qquad
 h\wedge f\longmapsto-2f,
\]
and is an isomorphism.  The next differential vanishes, since
\[
 d(e\wedge h\wedge f)
 =-2e\wedge f-h\wedge h-2f\wedge e=0.
\]
Together with $d:\mathfrak{sl}_2\to\mathbb C$ equal to zero, this gives
Chevalley homology dimensions $1,0,0,1$.  The ordinary bar--Chevalley
comparison identifies these homology groups with those of $U(\mathfrak{sl}_2)$;
it does not identify the two chain spaces.

\end{calculation}

\section{Positive Witt homology}
The positive Witt algebra is
\[
 L_1=\bigoplus_{n\ge1}\mathbb C e_n,
 \qquad[e_i,e_j]=(j-i)e_{i+j}.
\]
It is an ordinary positively graded Lie algebra with finite-dimensional weight pieces.
These are the positive modes of the centreless Virasoro bracket.
Its enveloping bar is an explicit ordinary associative construction;
identifying it with the transverse algebra of a particular chiral theory is
an additional comparison.

\begin{theorem}[Pentagonal support]
\label{thm:pentagonal-spine}
For $n\ge1$, the degree-$n$ homology $H_n(\BarPerp U(L_1))$ is one-dimensional in each
of the weights $(3n^2-n)/2$ and $(3n^2+n)/2$, and zero at every other weight.
Also $H_0=\mathbb C$ in weight zero. Its Euler character is
\begin{equation}\label{eq:witt-euler-spine}
 \chi_q\bigl(\BarPerp U(L_1)\bigr)=\prod_{m\ge1}(1-q^m).
\end{equation}
\end{theorem}
\begin{proof}
Theorem~\ref{thm:bar-chevalley-spine} identifies the bar homology with
Chevalley homology. In each weight the Chevalley complex is finite-dimensional,
so its homology is dual to the continuous cohomology in that weight.
Goncharova's main theorem and Example~1.1 compute the whole cohomology:
in degree $n$ there are precisely the two basic sequences
$(2,5,\ldots,3n-1)$ and $(3,6,\ldots,3n)$ in her exponent indexing
\cite{Goncharova1973}. Subtracting one from each exponent gives total weights
$(3n^2-n)/2$ and $(3n^2+n)/2$.
This proves both the two nonzero dimensions and the vanishing in all other weights.

Independently, the exterior Chevalley chain complex has Euler product
$\prod_{m\ge1}(1-q^m)$, because there is one even Lie generator in every
positive weight. Its coefficients are finite sums at each weight.
Euler's pentagonal identity expands it as
$\sum_{r\in\mathbb Z}(-1)^r q^{r(3r-1)/2}$, agreeing with the computed homology.
This agreement checks the Euler character; the vanishing assertion uses
Goncharova's homology theorem, not the Euler identity alone.
\end{proof}

For $n=1,2,3,4$ the weight pairs are $(1,2)$, $(5,7)$, $(12,15)$ and $(22,26)$.
The corresponding nonzero Euler coefficients at
$0,1,2,5,7,12,15,22,26$ have signs $+,-,-,+,+,-,-,+,+$.

\section{Planar path sequences}

A Motzkin path is a lattice path from $(0,0)$ to $(n,0)$, staying at
nonnegative height, with steps $(1,1)$, $(1,0)$, and $(1,-1)$.
Let $M(n)$ count all such paths of length $n$, and let $R(n)$ count those with
no horizontal step at height zero.  Both counts include the empty path:
$M(0)=R(0)=1$, while $M(1)=1$ and $R(1)=0$.
For $n\ge1$ consider the two sequences
\[
 a_n=M(n+1)-M(n)=1,2,5,12,30,76,\ldots,
 \qquad
 b_n=R(n+3)=3,6,15,36,91,\ldots.
\]

\begin{proposition}[Motzkin and Riordan identities]
\label{prop:one-sequence-spine}
For every $n\ge0$,
\[
 M(n)=R(n)+R(n+1),\qquad
 M(n+1)-M(n)=R(n+2)-R(n).
\]
The ordinary generating functions of the two path counts, and hence those
of $a_n$ and $b_n$, belong to
$\mathbb Q(z,\sqrt{1-2z-3z^2})$.
\end{proposition}
\begin{proof}
Work in the formal power-series ring $\mathbb Q[[z]]$, and write
$M(z)=\sum_{n\ge0}M(n)z^n$ and $R(z)=\sum_{n\ge0}R(n)z^n$.
A nonempty Motzkin path either starts with a horizontal step followed by
an arbitrary Motzkin path, or has a unique first-return decomposition
$U P D Q$, where $U$ and $D$ are the up and down steps,
$P$ and $Q$ are arbitrary Motzkin paths, and $P$ is drawn
one unit above the horizontal axis.  Thus
\[
 M=1+zM+z^2M^2.
\]
A nonempty path counted by $R$ has no initial horizontal step.  Its first
return similarly has the form $U P D Q$, with $P$ arbitrary and $Q$ counted
by $R$.  Consequently
\[
 R=1+z^2MR,\qquad R=(1-z^2M)^{-1}.
\]
The inverses here exist because their denominators have constant term one.
Using $z^2M^2=M-1-zM$, compute
\[
 (1+zM)(1-z^2M)
 =1+zM-z^2M-z^3M^2=1+z.
\]
It follows that $(1+z)R=1+zM$.  Comparing coefficients of $z^{n+1}$ gives
$R(n+1)+R(n)=M(n)$ for all $n\ge0$.  Subtracting consecutive identities
gives the difference formula.

Solving the quadratic equation for the formal series with constant term
one gives
\[
 M(z)=\frac{1-z-\sqrt{1-2z-3z^2}}{2z^2},\qquad
 R(z)=\frac{1+zM(z)}{1+z},
\]
where the square root has constant term one.  The apparent pole in the
formula for $M$ is removable as a formal series.  Finite index shifts and
differences preserve the indicated quadratic function field, proving the
last assertion.  These identities are consequences solely of the path
decompositions.
\end{proof}

\begin{remark}[Low-degree homology]
Theorem~\ref{thm:pentagonal-spine} gives homology dimensions $2,2,2$ in degrees $1,2,3$, whereas $a_1,a_2,a_3$ are $1,2,5$. These triples agree only in degree two. The degree-one discrepancy already excludes equality of the full homology dimension sequences.
\end{remark}

\begin{theorem}[Enveloping bars and planar path counts]
\label{thm:carrier-spine}
Let $B_\bullet(A)$ be the ordinary normalized bar over $\mathbb C$, with
trivial coefficients and homological degree given by bar length.  With
$a_n=M(n+1)-M(n)$ and $b_n=R(n+3)$ indexed by $n\ge1$ as above,
\[
 \dim H_1 B_\bullet(U(\mathfrak{sl}_2))=0\ne b_1=3,
 \qquad
 \dim H_1 B_\bullet(U(L_1))=2\ne a_1=1,
\]
where $L_1=\bigoplus_{m\ge1}\mathbb C e_m$ and
$[e_i,e_j]=(j-i)e_{i+j}$.
Hence neither sequence is the bar homology dimension sequence of the
corresponding enveloping algebra.  Neither is its full bar chain dimension
sequence: in both cases $B_n$ is countably infinite-dimensional for every
$n\ge1$.
\end{theorem}
\begin{proof}
For any augmented algebra $A$ with augmentation ideal $I$, the bar
differential out of $B_1=I$ is zero and the image from $B_2=I\otimes I$ is
$I^2$.  Therefore
\[
 H_1B_\bullet(A)=I/I^2.
\]
For an ordinary Lie algebra $\mathfrak g$, the inclusion of generators
induces an isomorphism
\[
 \mathfrak g/[\mathfrak g,\mathfrak g]\ \xrightarrow{\ \sim\ }
 I/I^2,\qquad I=\ker(U(\mathfrak g)\to\mathbb C).
\]
Indeed the generators span $I/I^2$, and
$[x,y]=xy-yx$ vanishes there.  For injectivity, equip
$\mathbb C\oplus\mathfrak g/[\mathfrak g,\mathfrak g]$ with the square-zero
product on its second summand.  The quotient map from $\mathfrak g$ to that
summand is a Lie map into the commutator Lie algebra of this associative
algebra.  The universal property of $U(\mathfrak g)$ extends it to an
augmented algebra map, which kills $I^2$ and provides an inverse on the
displayed quotient.

The brackets $[h,e]=2e$, $[h,f]=-2f$, and $[e,f]=h$ show that
$\mathfrak{sl}_2$ equals its commutator algebra.  For $L_1$, every nonzero
bracket has weight at least three.  Conversely, for $m\ge3$,
\[
 e_m=\frac{1}{m-2}[e_1,e_{m-1}].
\]
Thus $[L_1,L_1]=\bigoplus_{m\ge3}\mathbb C e_m$, and its abelianization
has basis the classes of $e_1,e_2$, in weights one and two.  This proves
both degree-one assertions without using any higher-degree calculation.

PBW gives a countably infinite basis of the augmentation ideal of either
enveloping algebra.  In the $L_1$ case the basis consists of nonempty
ordered monomials in the countable set $e_1,e_2,\ldots$, each monomial
having finite length.  Finite tensor powers of either augmentation ideal
are therefore countably infinite-dimensional.  The path counts are finite
in every degree, so they cannot count these full chain spaces.
\end{proof}

In particular, a complex whose homology dimensions are $b_n$ cannot be
quasi-isomorphic to $B_\bullet(U(\mathfrak{sl}_2))$, and one whose homology
dimensions are $a_n$ cannot be quasi-isomorphic to $B_\bullet(U(L_1))$,
with the stated degree convention.  A discrepancy of chain dimensions
alone would not imply this conclusion; the degree-one homology discrepancy
does.

The weight grading gives a different, finite comparison.  An ordered PBW
monomial of weight $w$ specifies a partition of $w$, so for $n\ge1$,
\[
 \sum_{w\ge0}\dim B_n(U(L_1))_w q^w
 =\left(\prod_{m\ge1}(1-q^m)^{-1}-1\right)^n.
\]
Each coefficient is finite.  Thus a weightwise dimension, a dimension
after a specified weight cutoff, and an ungraded dimension are distinct
invariants.  In a weight completion, dimensions must likewise be interpreted
weightwise; the countable algebraic dimensions above refer to the ordinary
direct-sum enveloping algebras.

\begin{remark}[Chiral coefficients and transverse multiplication]
The central charge of a Virasoro coefficient specifies no additional
transverse multiplication. The free, square-zero, enveloping and quantum-plane
examples give ordinary associative bars with different homology.
Realizing one beside a fixed chiral coefficient requires an actual coefficient
family and compatible operations; tensoring by an arbitrary fixed algebra
does not itself preserve factorization.
A comparison with an ordered geometric bar must construct that family and
a quasi-isomorphism preserving its transverse multiplication.
\end{remark}

\chapter{Information retained by comparison maps}
\label{ch:obstructions}

Dimensions, Euler characters and anomaly coefficients each forget structure.
The following examples identify that lost information and the extra maps
needed to recover it.
The ordinary algebra examples in this chapter are over $k=\mathbb C$.

\section{Alignment and exchange}
\begin{proposition}[The aligned retraction]
\label{prop:lambda-sector-spine}
In the finite decomposition model, let
$\iota:\mathcal M^{\mathrm{al}}\to\mathcal M$ and
$\pi:\mathcal M\to\mathcal M^{\mathrm{al}}$ be the inclusion and projection.
Then $\pi\iota=\id$ and $\iota\pi=e_{\mathrm{al}}$.
These maps are inverse on the full object precisely when its complementary
summand vanishes. A comparison on the aligned image alone does not specify
one on the complement or prove collision descent.
\end{proposition}
\begin{proof}
On an aligned summand both composites are the identity.
On an off-aligned summand $\iota\pi$ is zero.
This proves the identities and the criterion for this retraction to be an equivalence.
To obtain a mixed distributive law, one must separately verify its unit and
iteration diagrams and its compatibility with all coefficient transitions.
The retraction identities make no assertion excluding some other comparison
on the full object.
\end{proof}

ZF exchange gives concrete coherent crossings on a specified labelled sector.
To make it a mixed Ran operation one still needs the coefficient diagram and
its collision maps; the spectral word relation does not define those maps.

\begin{theorem}[Dimensions do not determine exchange]
\label{nogo:dimension-loss-spine}
The dimensions of the kernels of $V^{\otimes n}\to\Sym^nV$ do not determine
an $R$-operator, an associative multiplication or a Hopf structure on an algebra
built from $V$.
\end{theorem}
\begin{proof}
Take $V=kx\oplus ky$. The polynomial algebra $k[x,y]$ and the quantum plane
$A_q=k\langle x,y\rangle/(yx-qxy)$ have the same graded vector-space dimensions,
namely $\dim(A_q)_n=n+1$, by the PBW basis above.
For $q\ne1$, one is commutative and the other is not.
If $q$ has infinite multiplicative order, the centre of $A_q$ is $k$:
a monomial commuting with $x$ must have zero $y$-exponent, and commuting with
$y$ then forces zero $x$-exponent. The two module categories therefore also
have different centres.

On $V^{\otimes2}$ both the identity and the diagonal operator with
$q_{12}=q$, $q_{21}=q^{-1}$, $q_{11}=q_{22}=1$ satisfy the exchange equations.
The kernel dimensions depend only on $V$, so cannot select between them.
A coproduct or associator requires still further maps and coherence equations.
\end{proof}

\section{Coefficient operations and the centre}
\begin{proposition}[Topology and coefficient contractions]
The cellular equation $\partial^2=0$ does not specify operation-valued maps
on the strata of a compactified moduli space. Such maps define a differential
only when their signed composites satisfy the corner equation.
\end{proposition}
\begin{proof}
With constant coefficients the two incident face maps agree and their
orientation signs cancel. For coefficient maps the terms are their actual
composites; the calculation in the corner proposition gives the necessary
and sufficient equation. Their equality is extra data.
\end{proof}

\begin{proposition}[A boundary centre need not recover a bulk]
The derived centre represents central actions on a specified boundary algebra.
Its universal property does not imply that an acting algebra maps to it by
an equivalence.
\end{proposition}
\begin{proof}
For the semisimple algebra $B=k\times k$, take the boundary module $b=(k,0)$.
Its derived endomorphism algebra is $k$, and the action of the centre
$HH^0(B)=k\times k$ on it is projection to the first factor.
The second factor acts trivially. The same boundary algebra occurs for $B=k$,
whose centre is $k$. Thus the boundary algebra alone does not recover the
larger centre; $b$ is not a generator of $B$-modules.
A physical bulk-to-centre identification likewise needs its own equivalence theorem.
\end{proof}

\section{Central charge and anomaly classes}
\begin{proposition}[Equal central charge]
Central charge is not a complete invariant of a vertex operator algebra,
even among holomorphic theories.
\end{proposition}
\begin{proof}
The even unimodular rank-sixteen lattices $E_8\oplus E_8$ and $D_{16}^{+}$
give holomorphic lattice vertex operator algebras of central charge $16$
\cite{DongMasonSmall}. Their weight-one Lie algebras are respectively
$\mathfrak e_8\oplus\mathfrak e_8$ and $\mathfrak{so}_{32}$.
These Lie algebras are not isomorphic, so neither are the vertex operator algebras.
The same scalar anomaly therefore does not determine the chiral multiplication.
\end{proof}

A representation category, braid transport and conformal blocks require the
vertex operations and their representation data. The central charge specifies
the coefficient of the Virasoro central extension; it supplies none of those maps.

\begin{proposition}[The types of anomaly data]
\label{nogo:anomalies-spine}
A current-algebra level, Virasoro central charge, determinant-line curvature,
BRST obstruction and automorphic weight can be compared only after their
coefficient spaces and a comparison map have been specified.
Numerical agreement alone identifies none of these structured data.
\end{proposition}
\begin{proof}
A level scales a current cocycle, and a central charge scales the Virasoro cocycle.
A determinant curvature is a two-form representing the line's characteristic class.
A BRST obstruction lies in the cohomology of the deformation complex of the
BRST differential. An automorphic weight is a character of the relevant
structure group. These are different types of objects, not components of one
canonically given cohomology class. A construction relating two of them must
supply the intervening functor or cochain map and identify its image.
\end{proof}

\section{Bars and scalar sequences}
\begin{proposition}[A comparison of bar functors]
An identification of $\BarPerp(A^\perp)$ with
$\BarAssCh(A^{\ch}_{\mathsf{Ass}})$ requires a map relating their coefficient
categories, tensor products, augmentations and multiplication operations.
Agreement of Euler characters does not give such a map.
\end{proposition}
\begin{proof}
Both normalized constructions can have an ordinary word-length grading,
but their word factors and adjacent-face maps use different tensors and products.
A chain comparison must commute with those face maps and with the units;
a coalgebra comparison must also preserve the cut maps.
The Euler character retains only an alternating sum of dimensions.
\end{proof}

\begin{theorem}[A scalar sequence is not homology data]
\label{nogo:scalar-spine}
A chain complex, its differential and its grading must be specified before
an integer sequence can be identified with bar homology dimensions.
Equal chain dimensions, discriminants or growth rates do not determine homology.
\end{theorem}
\begin{proof}
The two complexes $[k\xrightarrow{0}k]$ and $[k\xrightarrow{1}k]$ have the
same chain dimensions and different homology.
For $U(\mathfrak{sl}_2)$, Calculation~\ref{calc:simple-spine} gives
countably infinite-dimensional full bar chain spaces in positive length.
The numbers $3^n$ count only their least total PBW-degree pieces.
Its Chevalley chain dimensions are $1,3,3,1$, while its homology dimensions
are $1,0,0,1$. The quasi-isomorphism identifies homology and not chain spaces.
Proposition~\ref{prop:one-sequence-spine} gives path sequences related by an
explicit transform; that transform is not a comparison with either bar complex.
\end{proof}

These examples leave positive constructions intact: the CE comparison,
path recursions, exchange normal forms and finite cotangent model each have
an explicit carrier and proof. A further geometric identification must construct
a map preserving the additional operations it claims to recover.
